\documentclass[11pt]{amsart}
\usepackage[T1]{fontenc}
\usepackage{amsmath,amssymb,mathrsfs,mathtools,booktabs,longtable,microtype}
\usepackage[margin=1.08in]{geometry}
\usepackage[hidelinks]{hyperref}
\hypersetup{pdftitle={General Theta Foundations I: Acquired Geometry and Causal Resource Transfer},pdfauthor={Qian Qi}}
\pdfinfoomitdate=1
\pdftrailerid{}
\pdfsuppressptexinfo=15
\newtheorem{theorem}{Theorem}[section]
\newtheorem{proposition}[theorem]{Proposition}
\newtheorem{lemma}[theorem]{Lemma}
\newtheorem{corollary}[theorem]{Corollary}
\theoremstyle{definition}
\newtheorem{definition}[theorem]{Definition}
\newtheorem{example}[theorem]{Example}
\theoremstyle{remark}
\newtheorem{remark}[theorem]{Remark}
\newcommand{\E}{\mathbb E}
\newcommand{\PP}{\mathbb P}
\newcommand{\R}{\mathbb R}
\newcommand{\N}{\mathbb N}
\newcommand{\cF}{\mathcal F}
\newcommand{\cH}{\mathcal H}
\newcommand{\TV}{\operatorname{TV}}
\newcommand{\law}{\operatorname{Law}}
\newcommand{\osc}{\operatorname{osc}}
\newcommand{\diam}{\operatorname{diam}}
\newcommand{\dist}{\operatorname{dist}}
\newcommand{\cp}{\mathrm{cp}}
\newcommand{\on}{\mathrm{on}}
\newcommand{\tr}{\mathrm{tr}}
\newcommand{\one}{\mathbf 1}
\newcommand{\norm}[1]{\left\lVert#1\right\rVert}
\newcommand{\abs}[1]{\left\lvert#1\right\rvert}
\title{General Theta Foundations I:\\Acquired Geometry and Causal Resource Transfer}
\author{Qian Qi}
\date{7 October 2026}
\subjclass[2020]{Primary 60G35, 62B05; Secondary 94A34, 93E11, 28A80}
\keywords{causal experiment, predictive quotient, acquired geometry, recursive quantization, singular acquisition, resource transfer}
\begin{document}
\begin{abstract}
We study the resources required to retain predictions acquired by a prepared causal experiment. A block transfer theorem constructs a hard finite-state encoder from static covers and inequalities for the physical report kernels. The construction quantizes after every raw call; no exact posterior or observation window is retained between calls. A relation preserved by the cover controls the approximate trajectory through blocks whose individual updates need not contract. Actual acquired mass gives matching checkpoint converses, while conditional projection separates acquisition uncertainty from memory distortion, including finite atomic acquisition laws. We verify the theorem for sparse finite-state hidden Markov experiments with Gaussian observations and state-dependent missingness, and for finite-read singular experiments with expanding deletion and nonreset transport. A stopped drift bound records the cost of adaptive termination. Parameter-independent causal transducers, their clocks and task maps are fully specified; complete stopped-law defects compose in the common decision metric. Matching statements concern declared task, memory, acquisition and output interfaces; all other computational resources are charged separately.
\end{abstract}
\maketitle
\section{The problem and the experimental object}\label{sec:experiments}
The geometry relevant to prediction is the geometry of information that an experiment has actually acquired. It is not, in general, the geometry of the parameter space, of all formally reachable posteriors, or of an arbitrarily selected reference measure. A second distinction is equally important. A finite list of accurate terminal predictions does not itself give a causal implementation of that list: the reports needed to select a representative may already have been discarded.

Our main result, Theorem~\ref{thm:main}, connects these two issues. It is a sufficient raw-kernel construction from a static global cover, block two-point stability, and a cover-preserved relation which admits a robust envelope bound. Unlike a one-step contraction theorem, it allows every individual transition in a block to have contraction coefficient one. Unlike a blockwise coding prescription, its encoder retains only one finite label after \emph{each} physical report. The proof compares its perturbed trajectory with an exact shadow started at the beginning of a block; a relation on candidate states, independent of the label budget, controls the perturbations inside the block. Actual acquired small-ball mass then gives the converse in the same squared task.

The distinction has a concrete consequence. Sparse transition matrices need not mix in one step, and continuous Gaussian observations have unbounded log-likelihood ratios. Nevertheless, a primitive support pattern and a positive-probability bounded observation block suffice. A mass-preserving inward compander prevents intermediate rounding from destroying the positivity accumulated over that block. The resulting hard-label law has exponent determined by a positive submeasure of the actual posterior distribution. Theorem~\ref{thm:hmm} proves this for an arbitrary fixed finite dimension, with time-varying sparse matrices and state-dependent missing observations. Theorem~\ref{thm:singular} verifies the same general mechanism on an entirely different, singular predictive domain. There the actual acquired law is finite atomic; a physical-oracle comparison, not a substitution of a Cantor law for that atomic law, supplies the geometric converse.

The general theorem is followed by exact acquisition identities, a stopped drift estimate, and a composition theorem for fully specified causal transducers. Uniform deterministic-time risk is distinguished from stopped risk, and a complete report-law comparison is distinguished from a prediction bound. The structural hypotheses are not asserted to be necessary for every finite-memory predictable experiment. In particular, exact copying can preserve a code without any strict contraction. This boundary does not reduce the experimental problem; it identifies which mechanism the present theorem proves.

\subsection{Prepared kernels and legal strategies}
All measurable spaces below are standard Borel. A parameter-indexed experiment consists of hidden spaces $X_t$, action spaces $A_t$, report spaces $Y_{t+1}$, a preparation $\mu_\lambda$, and measurable probability kernels
\begin{equation}\label{eq:raw}
 K_t^\lambda(dx',dy,dc\mid x,a).
\end{equation}
The cost $c$ takes values in a declared finite product of nonnegative extended real spaces and includes elapsed physical time. Positivity means nonnegativity and normalization; no strict lower bound is imposed at this level. Failures, missing observations and Dirac transitions are legitimate reports. A visible history records actions, reports and public costs. A continuation interface records available actions, remaining resources and any clock which changes the legal future.

A strategy is a measurable kernel $\beta_t(da\mid h_t)$ supported on the available actions. It may use independent randomization but cannot read $\lambda$, an unobserved hidden state, or a future report. A stopping rule is a stopping time for this visible filtration, with the declared stop action and terminal test. Kernel integration defines the actual path law. For the Bayesian prediction statements, fix a preparation (or a prior on $(\lambda,X_0)$) and a legal exploration rule \emph{independently of the encoder}. A posterior for this preparation is not automatically a parameter-uniform sufficient statistic. The latter quantifiers occur explicitly in Section~\ref{sec:morphisms}.

A persistent $M$-label encoder has measurable updates $I_t=q_t(I_{t-1},a_t,y_t)$ with at most $M$ possible labels at each cut. Time-dependent read-only programs are allowed only with their description and clock accounted for. Randomized updates are allowed unless stated otherwise. No earlier report can be reread. A checkpoint encoder inspects the actual terminal history once and retains at most $M$ labels. Neither model includes a free real posterior, unbounded random seed containing data, or hidden observation tape. A supplied controller's state is charged separately from the prediction label.

\subsection{Executable future tests and a measurable quotient}
An executable test is a legal finite continuation followed by a bounded function of its actual reports. Histories are equivalent when they have the same continuation interface and all conditional test expectations agree. This definition concerns an operational class of tests, not arbitrary functions of a hidden state. The next certificate supplies a measurable realization; applications may instead verify their quotient directly.

\begin{proposition}[Continuous-instrument quotient]\label{prop:quotient}
Let $\mathcal B$ be compact metrizable. Let $\mathcal V\subset C(\mathcal B)$ be a separable closed linear space generated by continuous executable-test evaluations and interface coordinates, containing constants. A countable dense family determines all declared tests. For each action $a$, let the report domain be Polish with reference measure $m_a$, and let $g_a(\pi,y)$ be a jointly continuous nonnegative density. On $D_a=\{g_a>0\}$ suppose the normalized belief update $B_a$ is continuous. For every $f\in\mathcal V$ and report $y$, require the instrument evaluation $g_a(\pi,y)f(B_a(\pi,y))$, extended by its instrument value at density zero, to be continuous in $\pi$ and in $\mathcal V$. Require joint continuity of these extended evaluations where continuity in the report is used. Then the evaluation image $S$ is compact metrizable, its fibers are predictive equivalence, and the report laws and normalized updates descend to $S$. The descended update is continuous on the open positive-density domain in $S\times Y$.

If a standard-Borel statistic measurably determines every coordinate evaluation, it measurably determines $S$. The assertion may be imposed on one common full-measure history set for an almost-sure version.
\end{proposition}
\begin{proof}
Normalize the determining evaluations to lie in $[-1,1]$ and let $E:\mathcal B\to[-1,1]^{\N}$ be their product. Its image is compact and hence Borel. Agreement on its coordinates implies agreement on $\mathcal V$ by uniform approximation and on the declared tests by the determining hypothesis. Taking $f=1$ shows that equal fibers have the same density at every report. Dividing the equal instrument evaluations by that positive common density shows that $E(B_a(\pi,y))$ is constant on the corresponding fibers.

The continuous surjection $E$ is closed, with compact fibers. The product $E\times\mathrm{id}_Y$ is a closed quotient map: given a closed set and a convergent image point, compactness of the belief fiber supplies a convergent subnet witnessing closedness locally. Equivalently one can use the finite subcover characterization of a closed map with compact fibers. The joint instrument coordinates therefore descend continuously. Their ratios by the positive descended density are continuous on the open set where it is positive; these ratios are precisely the coordinates of the descended update. No claim is made about normalization at a null report. A jointly Borel null-report extension, when needed by a numerical representative, is additional data.

For minimality let $u_n$ be measurable coordinate factors through the given statistic. The product $u=(u_n)$ is measurable. Replace $u$ by a fixed state outside the Borel set $S$ to obtain an $S$-valued factor. Countably many almost-sure equalities can be imposed on one common full-measure set. Thus the conclusion does not rely on measurability of an unspecified set-theoretic quotient.
\end{proof}

The metric used for stability need only generate the relevant Borel structure, not the compact quotient topology. Log-ratio distance on the open probability simplex is an important example. Report versions used to update a wrong representative must be defined jointly, rather than selected separately almost surely for each state.

\subsection{One scored task}
Let $Y$ be a bounded scalar terminal report, or let an independently selected finite probe index specify one of finitely many bounded scalar reports. In the latter case use the weighted Euclidean norm of their separately executable squared losses. No joint counterfactual measurement is assumed. Put
\begin{equation}\label{eq:task}
 \begin{gathered}
 P_T=\E[Y\mid\cH_T],\quad \nu_T=\law(P_T),\quad B_T=\E\norm{Y-P_T}^2,\\
 q_M(\nu)=\inf_{\#C\le M}\int\dist(p,C)^2\nu(dp).
 \end{gathered}
\end{equation}
The infimal checkpoint and online risks are $R_\cp(T,M)$ and $R_\on(T,M)$. The law $\nu_T$ is always the pushforward of the actual preparation and exploration. It may be atomic, mixed, singular, transient, or supported on several strata. Its role in the theorem is probabilistic rather than formal-dimensional.

\section{A causal acquired-geometry transfer theorem}\label{sec:main}
Let $(S,d)$ be a separable metric predictive domain with its standard Borel structure. Conditional on the complete observed past and used algorithm randomization, the physical report at step $t$ has law $J_t(dy\mid S_{t-1},a_t)$ and the predictive update is $F_t(S_{t-1},a_t,y)$. Both objects are jointly Borel on the entire declared true and candidate domain. All comparisons below use the \emph{physical} report law and the same realized actions. They are uniform over the declared strategy class, including its legal history-dependent choices. The action sequence in a shadow flow is the one selected from the physical reports, not a newly optimized shadow policy.

Fix a block length $m\ge1$, a Borel envelope $w:S\to[1,\infty)$, and a Borel relation $\mathcal D\subset S\times S$, writing $z\in\mathcal D(v)$ for $(v,z)\in\mathcal D$. Require $v\in\mathcal D(v)$ and $w(z)\le w(v)$ for $z\in\mathcal D(v)$. For each budget, static data comprise at most $M$ representatives $c_i$ and Borel cells $A_i$ covering $S$, with
\begin{equation}\label{eq:static}
 v\in A_i\quad\Longrightarrow\quad
 c_i\in\mathcal D(v),\qquad d(v,c_i)\le r_M w(v).
\end{equation}
This is not a recursive machine assumption. The relation is fixed independently of $M$ and of the desired error radius.

Write $\mathscr F_{s,j}(z)$ for the exact composition of $j$ updates starting at $z$ at time $s$, along the next physical actions and reports. The hypotheses are the following explicit kernel inequalities. They hold conditionally on every declared past at block starts $s=km$, for every true state $x$ and candidate $z$ compatible with that past.

First, for some $L\ge1$ each one-step map is pathwise $L$-Lipschitz on the common report domain, and
\begin{equation}\label{eq:blockpair}
 \E_x[d(\mathscr F_{s,m}(x),\mathscr F_{s,m}(z))^2\mid\cF_s]
 \le\kappa d(x,z)^2,\qquad 0\le\kappa<1.
\end{equation}
The expectation includes the actual report-driven action choices. Second, consider \emph{any} adapted selections, starting at $z_0=z$,
\[
 v_j=F_{s+j}(z_{j-1},a_{s+j},Y_{s+j}),\qquad
 z_j\in\mathcal D(v_j),\quad 1\le j\le m.
\]
For constants $0\le a<1$, $b\ge0$, $A,B\ge0$, require
\begin{align}
 \E_x[w(z_m)^2\mid\cF_s]&\le a w(z)^2+b,\label{eq:robustdrift}\\
 \E_x[w(v_j)^2\mid\cF_s]&\le A w(z)^2+B,
                 \qquad 1\le j\le m.\label{eq:within}
\end{align}
These inequalities quantify over arbitrary relation-preserving perturbations, not over the error of a supplied encoder. They are verifiable from raw positivity or envelope estimates. Their strength and their verification cannot be suppressed: ordinary inwardness by itself need not preserve the mixing accumulated across several sparse transitions.

A finite-input implementation may replace exact selection by a declared Borel routine with values in the same codebook, still in $\mathcal D(v_j)$, and satisfying
\begin{equation}\label{eq:localerror}
 d(v_j,z_j)\le r_M w(v_j)+u.
\end{equation}
Here $u$ is a certified joint bound for numerical, input and calibration errors in the stable metric. Borel exact selection has $u=0$. Finite workspace and computable membership are additional implementation certificates, not consequences of measurability.

\begin{theorem}[Block construction and acquired-risk transfer]\label{thm:main}
Fix the experiment, a deterministic horizon $T$, and exploration independent of the encoder. Suppose \eqref{eq:static}--\eqref{eq:within} hold and a legal seed in the codebook has
\[
 \E w(\widehat S_0)^2\le H_0,\qquad
 \norm{d(S_0,\widehat S_0)}_{L^2}\le e_0.
\]
All seed acquisitions and retained states are charged. Let $P_T=f_T(S_T)$, with $f_T$ $L_f$-Lipschitz in the task norm. Put
\begin{align}
 H&=\max\{H_0,b/(1-a)\},&H_*&=AH+B,\nonumber\\
 \Lambda_j&=\sum_{i=0}^{j-1}L^i\quad(\Lambda_0=0),&
 \Delta&=r_M\sqrt{H_*}+u,\nonumber\\
 E_*&=\max\{e_0,\Lambda_m\Delta/(1-\sqrt\kappa)\},&
 \mathcal E_M&=L^{m-1}E_*+\Lambda_{m-1}\Delta.\label{eq:constants}
\end{align}
The static cover constructs an $M$-label causal encoder, quantized after every raw call, for which
\begin{equation}\label{eq:mainrisk}
 B_T+q_M(\nu_T)=R_\cp(T,M)
 \le R_\on(T,M)
 \le B_T+(L_f\mathcal E_M+v)^2,
\end{equation}
where $v$ is an additional root-mean-square terminal readout error. The constants are uniform in $T$ and in the certified strategies, not uniform as $\kappa$ or $a$ tend to one.

If an actual submeasure $\underline\nu_T\le\nu_T$ has mass $\zeta_T>0$ and
\begin{equation}\label{eq:smallball}
 \sup_c\underline\nu_T(B(c,s_M))\le\zeta_T/(2M),
\end{equation}
then $R_\cp(T,M)-B_T\ge\zeta_Ts_M^2/2$. Thus $e_0=O(r_M)$, $u,v=O(r_M)$, $s_M\gtrsim r_M$, and $\inf_T\zeta_T>0$ give matching checkpoint and online excess risk of order $r_M^2$. Every constant is determined by the declared experiment, cover and actual acquired witness.
\end{theorem}

\begin{lemma}[Checkpoint projection]\label{lem:projection}
For every forecast based on the history and randomization independent of the experiment and exploration,
\begin{equation}\label{eq:projection}
 \E\norm{Y-a}^2=B_T+\E\norm{P_T-a}^2.
\end{equation}
The optimal checkpoint distortion is $q_M(\nu_T)$, including independent shared randomization.
\end{lemma}
\begin{proof}
Condition on the history and the algorithm randomization; the conditional mean remains $P_T$, so the cross term vanishes. For each value of an independent shared seed, replace a randomized decoder by its label-conditional mean. There are at most $M$ centers, and random label selection cannot outperform the closest center. A first-nearest-center selector is Borel. Every center list therefore gives distortion at least $q_M(\nu_T)$, and approximating lists give the reverse inequality. Independence keeps $\nu_T$ fixed when conditioning on this seed. Randomness used to choose the exploration is not a free such seed; any disclosed exploration side-information must be conditioned on or charged as an interface.
\end{proof}

\begin{lemma}[Block shadow and accumulated rounding]\label{lem:block}
Let $Q_M$ choose the first cell in \eqref{eq:static}, and set
\begin{equation}\label{eq:encoder}
 \widehat S_t=Q_M(F_t(\widehat S_{t-1},a_t,Y_t)).
\end{equation}
The same assertions hold for the certified routine \eqref{eq:localerror}. At block endpoints $\E w(\widehat S_{km})^2\le H$. All pre-rounding candidates have second envelope moment at most $H_*$. If $e_t=\norm{d(S_t,\widehat S_t)}_2$, then
\begin{align}
 e_{s+m}&\le\sqrt\kappa\,e_s+\Lambda_m\Delta,\label{eq:recurrence}\\
 e_{s+j}&\le L^j e_s+\Lambda_j\Delta,\qquad 0\le j<m.\label{eq:interior}
\end{align}
\end{lemma}
\begin{proof}
First-index selection from a finite Borel cover is Borel. It obeys the fixed relation, so \eqref{eq:robustdrift} applies to its trajectory and gives the endpoint envelope induction. Equation~\eqref{eq:within} gives the pre-rounding bound at every microstep.

Start an exact shadow at $\widehat S_s$ and feed it the \emph{same} physical actions and reports for $m$ steps, without storing that shadow in the implementation. Successively replace the perturbed transitions by exact ones, from the last insertion of error backwards. Pathwise Lipschitz continuity gives
\[
 d(\mathscr F_{s,m}(\widehat S_s),\widehat S_{s+m})
 \le\sum_{j=1}^m L^{m-j}
   d(F_{s+j}(\widehat S_{s+j-1},a_{s+j},Y_{s+j}),\widehat S_{s+j}).
\]
This is a telescoping comparison of functions of the observed word, not a claim that reports are independent of the error. Each summand has $L^2$ norm at most $\Delta$ by \eqref{eq:localerror} and the envelope bound. Minkowski and \eqref{eq:blockpair} give \eqref{eq:recurrence}. Applying the same telescoping argument for $j<m$, with the pathwise gain $L^j$ for the initial discrepancy, gives \eqref{eq:interior}. The exact shadow is only a proof device; no shadow state or block of reports is retained by \eqref{eq:encoder}.
\end{proof}

\begin{proof}[Proof of Theorem~\ref{thm:main}]
The recurrence preserves $e_{km}\le E_*$ by the definition of $E_*$. Equation~\eqref{eq:interior} gives $e_T\le\mathcal E_M$. The readout $f_T(\widehat S_T)$ has root task error at most $L_f\mathcal E_M$, and its certified terminal perturbation adds $v$ in that same norm. Lemma~\ref{lem:projection} proves the upper and exact checkpoint identity. Every online output is an allowed checkpoint output, which gives the middle inequality.

For any $M$ centers, their open $s_M$-balls cover at most half the submeasure in \eqref{eq:smallball}. Outside their union the squared distance to the list is at least $s_M^2$. Its contribution is at least $\zeta_Ts_M^2/2$. Seed conditioning as in Lemma~\ref{lem:projection} proves the randomized converse. This argument uses no volume measure or manifold beyond the actual submeasure specified in the hypothesis.
\end{proof}

\begin{corollary}[Exact independent suspensions]\label{cor:holds}
For a one-step certificate suppose active reports satisfy a pair bound $\kappa<1$ and a pre-rounding envelope bound $aw(z)^2+b$. At physical step $t$ an active indicator, independent of the complete past, has deterministic probability $\vartheta_t\in[0,1]$; on inactivity the predictive and retained states are copied exactly. Then
\[
 H=\max(H_0,b/(1-a)),\qquad
 e_t\le\max\{e_0,(r_M\sqrt H+u)/(1-\sqrt\kappa)\}
\]
for every deterministic $t$, with all inactive calls and time charged.
\end{corollary}
\begin{proof}
Use indicator-form expectations. The candidate-envelope inequality gives
\[
\E[\one_{\{D_t=1\}}w(v_t)^2]
\le\vartheta_t(a\E w(\widehat S_{t-1})^2+b).
\]
On inactivity the old envelope is unchanged. This proves the invariant bound, including $\vartheta_t=0$ without conditioning on a null event. The squared error obeys
\[
 e_t^2\le(1-\vartheta_t)e_{t-1}^2+
 \vartheta_t(\sqrt\kappa e_{t-1}+r_M\sqrt H+u)^2.
\]
Both terms preserve the displayed radius. Independence from the past is used here and is not inferred from a state-dependent missing-report mechanism.
\end{proof}

For $m=1$ and $\mathcal D(v)=\{z:w(z)\le w(v)\}$, the raw one-step envelope bound verifies the relation-robust condition automatically. Thus the one-step inward theorem and its suspension mechanism are retained. The block theorem adds a new route when accumulated positivity, rather than one-step strict contraction, is the available experimental mechanism.

\section{Acquisition, finite outputs, and adaptive termination}\label{sec:acquisition}
\subsection{Actual mass and physical-oracle decomposition}
\begin{proposition}[Acquired geometry and acquisition uncertainty]\label{prop:acquired}
If $\underline\nu\le\nu$ has mass $\zeta>0$ and
$\sup_c\underline\nu(B(c,r))\le C r^d$ for $0<r\le r_0$, then
\begin{equation}\label{eq:masslower}
 q_M(\nu)\ge\frac\zeta2\left(\frac\zeta{2CM}\right)^{2/d}
\end{equation}
whenever the displayed radius is at most $r_0$.

More generally, let $\cH\subset\mathcal G$ be the actual acquisition and a declared physical-oracle sigma field. Put
\[
 X=\E[Y\mid\mathcal G],\quad P=\E[X\mid\cH],\quad
 b_\circ=\E\norm{Y-X}^2,\quad A=\E\norm{X-P}^2.
\]
Then
\begin{equation}\label{eq:oracle}
 R_\cp(M)-b_\circ=A+q_M(\law(P))\ge q_M(\law(X)).
\end{equation}
If the online upper is $A+C_1r_M^2$ and $q_M(\law(X))\ge c_1r_M^2$, checkpoint and online excess over $b_\circ$ both have order $A+r_M^2$. The acquisition law $\law(P)$ may be finite atomic.
\end{proposition}
\begin{proof}
The first conclusion is the union-of-balls argument with radius $(\zeta/(2CM))^{1/d}$. Conditional projection twice gives
\[
 \E\norm{Y-a}^2=b_\circ+A+\E\norm{P-a}^2
              =b_\circ+\E\norm{X-a}^2.
\]
An $M$-label readout supplies at most $M$ centers also for $X$, proving \eqref{eq:oracle}. The left side is at least both $A$ and $c_1r_M^2$, hence at least a fixed multiple of their sum. The online upper supplies the reverse comparison. The ideal law is used only as a lower witness, never as the actual predictive law.
\end{proof}

\begin{proposition}[Output and calibration cuts]\label{prop:calibration}
A fixed finite output set $G$, independent of the acquired history, changes the nominal checkpoint distortion to
\[
 q_{M,G}(\nu)=\min_{C\subset G,\ \#C\le M}\int\dist(p,C)^2\nu(dp)
 \ge q_{\min(M,\#G)}(\nu).
\]
An $h$-net output set adds at most $h$ to root task error.

For a scalar latent score $X$ with $\E X=1/2$ and $X\pm\delta\in[0,1]$, suppose acquisitions are independent of an unknown sign $s\in\{-1,1\}$ and the final report is Bernoulli with mean $X+s\delta$. Let $P=\E[X\mid\cH]$, $A=\E(X-P)^2$ and $b_\delta=\E[X(1-X)]-\delta^2$. For every sign-independent forecast,
\begin{equation}\label{eq:sign}
 \sup_s\E_s(Y-a)^2-b_\delta
 =A+\E(P-a)^2+\delta^2+2\delta\abs{\E(P-a)}.
\end{equation}
Thus the excess lies between $A+\E(P-a)^2+\delta^2$ and $A+2\E(P-a)^2+2\delta^2$. The same bounds hold for a finite weighted probe task with a common sign.
\end{proposition}
\begin{proof}
Expected loss is affine in each label's output distribution on $G$, so a deterministic output is optimal; nearest-center selection proves the first formula. Rounding gives the root-error bound by Minkowski. For the sign formula the physical-oracle variance is the same $b_\delta$ for either sign, because $\E X=1/2$. Projection and expansion give
$b_\delta+A+\E(P-a)^2+\delta^2+2s\delta\E(P-a)$.
Maximizing proves \eqref{eq:sign}. Cauchy--Schwarz and $2\delta e\le\delta^2+e^2$ prove the upper. For weighted probes, expand componentwise and apply Cauchy--Schwarz to the weighted mean. A nominal conditional-mean decoder has zero mean error and attains the lower sign expression. A biased grid decoder need not do so.
\end{proof}

A calibration \emph{allowance} alone does not force a $\delta^2$ lower bound. The inaccessible two-sign task is the specified witness which does. Likewise an output-grid cut is not a lower bound on every internal arithmetic operation.

\subsection{A stopped drift theorem}
Uniform bounds at deterministic times do not by themselves control a time selected from the observations. The next result retains the actual stopped acquisition and explicitly pays for that selection. The exploration, including its stop rule, is fixed independently of the prediction encoder. A stop action must leave the terminal forecast dependent only on the charged retained tuple and the executed probe index.

\begin{theorem}[Bounded adaptive termination]\label{thm:stopping}
Assume Theorem~\ref{thm:main} with a certificate valid up to $mK$ raw calls. Every allowed stopped exploration must be a prefix of a declared legal continuation to $mK$ for which that certificate holds; this continuation is used only in the proof. Define
\begin{align*}
 \rho&=r_M+u,& \bar\kappa&=(1+\kappa)/2,&
 C&=\bar\kappa/(\bar\kappa-\kappa),\\
 \gamma&=\max\{\bar\kappa,(1+a)/2\},&
 c&=\max\{1,C\Lambda_m^2A/(\gamma-a)\},\\
 c_0&=C\Lambda_m^2(B+1)+cb,&
 V_k&=d(S_{km},\widehat S_{km})^2+c\rho^2w(\widehat S_{km})^2.
\end{align*}
Then $\E[V_{k+1}\mid\cF_{km}]\le\gamma V_k+c_0\rho^2$. For every block-boundary stopping time $\sigma\le K$,
\begin{equation}\label{eq:stoppedblock}
 \E V_\sigma\le\E V_0+c_0\rho^2\E\sigma,
 \qquad
 (1-\gamma)\E\sum_{k<\sigma}V_k
 \le\E V_0+c_0\rho^2\E\sigma.
\end{equation}
For an arbitrary visible stopping time $\tau\le mK$, put $\sigma=\lceil\tau/m\rceil$. With
\[
 C_1=2m(L^{2m}+\Lambda_m^2A/c),\qquad
 C_2=2m\Lambda_m^2(B+1),
\]
one has
\begin{equation}\label{eq:rawstop}
 \E d(S_\tau,\widehat S_\tau)^2
 \le\left(1+\frac{C_1}{1-\gamma}\right)\E V_0
 +\left(\frac{C_1c_0}{1-\gamma}+C_2\right)\rho^2\E\sigma.
\end{equation}
Consequently a uniformly $L_f$-Lipschitz stopped task has risk at most $B_\tau+(L_f\sqrt{\text{right side of \eqref{eq:rawstop}}}+v)^2$. All calls up to $\tau$, failures, preparation, terminal measurement and the stopping/clock state are counted.
\end{theorem}
\begin{proof}
Conditionally on a block-start past, the proof of Lemma~\ref{lem:block} gives the root bound
\[
 \norm{d(S_{s+m},\widehat S_{s+m})}_{2\mid\cF_s}
 \le\sqrt\kappa\,d(S_s,\widehat S_s)
     +\Lambda_m\rho\sqrt{Aw(\widehat S_s)^2+B+1}.
\]
Indeed $r_M\sqrt{x}+u\le(r_M+u)\sqrt{x+1}$. The inequality
$(\sqrt\kappa d+z)^2\le\bar\kappa d^2+Cz^2$, also valid at $\kappa=0$, and \eqref{eq:robustdrift} give the claimed drift by the choice of $c$. Multiply its rearranged form by $\one_{\{k<\sigma\}}$, which is known at time $km$, and sum. Bounded stopping and finite moments justify all sums. Dropping either the nonnegative terminal term or the negative drift gives \eqref{eq:stoppedblock}.

For the raw stopping rule, $\{\sigma\le k\}=\{\tau\le km\}$, so $\sigma$ is a bounded stopping time for the block filtration. Within a block, telescoping and conditional Minkowski, followed by $(x+y)^2\le2x^2+2y^2$, show
\[
 \E\left[\max_{1\le j\le m}d(S_{km+j},\widehat S_{km+j})^2\mid\cF_{km}\right]
 \le C_1 V_k+C_2\rho^2.
\]
Here the maximum is bounded by the sum of the $m$ nonnegative squared errors. The event $\{\tau>km\}$ is known at the block start. Bound the error at $\tau>0$ by the sum of these block maxima over $k<\sigma$, and the error at $\tau=0$ by $V_0$. Apply the second inequality in \eqref{eq:stoppedblock} to obtain \eqref{eq:rawstop}. Projection at the stopped sigma field and the terminal perturbation give the risk statement. One can extend the experiment after a stop by any certified legal policy for the proof; only the actual pre-stop calls are charged or used in the output.
\end{proof}

The exact checkpoint identity remains $B_\tau+q_J(\nu_\tau)$ when \emph{all} terminal information available to the decoder has at most $J$ retained states. In particular a prediction label of size $M$ and a separately retained stopping clock of size $Q$ give a $J=MQ$ cut, not an unqualified $M$-center formula. When a declared public terminal interface is free, the identity must instead condition on that interface. The stopped law is the actual law induced by the fixed exploration; it is not the deterministic-time law evaluated at a random integer after the calculation.

\begin{example}[Why stopping requires its own estimate]\label{ex:stop}
Let each report $Z_t$ be an independent Bernoulli variable with success probability $r^2$, and consider the fixed forecast zero. Its error for the reported state at any deterministic time is $r^2$. Stop at the first success or at $\lceil r^{-2}\rceil$. Its stopped error is $1-(1-r^2)^{\lceil r^{-2}\rceil}$, bounded away from zero. This is a statement about transfer of an existing deterministic-time error estimate, not a lower bound against every alternative stopped predictor. It explains the explicit acquisition/selection term in \eqref{eq:rawstop}.
\end{example}

\section{Sparse filters with endogenous missing observations}\label{sec:filters}
We verify the block theorem on a standard hidden-state observation class, without a positive entry floor on the \emph{zero} entries of its transition matrices and without bounded likelihood ratios. Fix $n=d+1\ge2$. Let $P$ be a primitive zero--one support matrix, with $P^m$ strictly positive. At each time a declared action selects a stochastic matrix $Q_t$ with support $P$; each nonzero entry is at least $s>0$. Thus every realized product of $m$ such matrices has every entry at least $s^m$. The matrices may vary with time and with past reports; no common invariant law is assumed.

Prepare the hidden state $Z_0$ uniformly on $\{0,\ldots,d\}$. After its transition, a state $i$ is reported as missing, $Y_t=\bot$, with probability $1-q_i^{a_t}$, and otherwise produces a Gaussian report in $\R^d$ with mean $\mu_i^{a_t}$ and covariance $\sigma^2I_d$. All kernels are known. Assume
\begin{equation}\label{eq:hmmparams}
 0<q_-\le q_i^a\le q_+<1,\qquad
 \norm{\mu_i^a}_\infty\le R,
 \qquad \abs{\det B_a}\ge b_0>0,
\end{equation}
where row $i$ of $B_a$, $1\le i\le d$, is $(\mu_i^a-\mu_0^a)^{\mathsf T}/\sigma^2$. A finite action set is sufficient. The positive determinant hypothesis is used for a $d$-dimensional acquired lower bound, not for the upper construction. A fixed family of affinely independent Gaussian means verifies it; no preassigned Fisher or covariance geometry of the predictive quotient is assumed.

Relative to Lebesgue measure plus a unit atom at $\bot$, the emission density is
\begin{equation}\label{eq:emissions}
 g_i^a(y)=q_i^a(2\pi\sigma^2)^{-d/2}
 e^{-\norm{y-\mu_i^a}^2/(2\sigma^2)},\quad
 g_i^a(\bot)=1-q_i^a.
\end{equation}
Missingness therefore reveals information about the hidden state when the $q_i^a$ differ. It is not an independent exact hold. The preparation, every transition/report attempt and the terminal probe each cost one raw acquisition. For a horizon $T$ the raw count is $T+2$.

Let $\pi_{t,i}=\PP(Z_t=i\mid\cH_t)$ and let $\overline\pi_{T,i}=\PP(Z_T=i)$ under the fixed exploration. The posterior is positive, and Bayes' formula gives the pointwise update
\begin{equation}\label{eq:bayes}
 F_t(\pi,a,y)_i=\frac{(\pi Q_t)_i g_i^a(y)}{\sum_j(\pi Q_t)_j g_j^a(y)}.
\end{equation}
The time and action interface accompany the posterior. A terminal probe chooses $i\in\{1,\ldots,d\}$ uniformly and emits a Bernoulli report with mean
\begin{equation}\label{eq:filtertask}
 X_i+\theta,\qquad
 X_i=\tfrac12+\tfrac14(\one_{\{Z_T=i\}}-\overline\pi_{T,i}),
 \quad\theta\in\{-\delta,\delta\},\quad0\le\delta\le1/8.
\end{equation}
No acquisition reveals $\theta$. The posterior determines every continuation by kernel integration and is separated by these executable probes. It is therefore a direct realization of the predictive quotient. The task norm is $\norm{x}_d^2=d^{-1}\sum_i x_i^2$.

On the open simplex set
\[
 d_H(\pi,\rho)=\osc_i\log(\pi_i/\rho_i),\qquad
 D(\pi)=\osc_i\log\pi_i,\qquad w(\pi)=e^{\eta D(\pi)}.
\]
This is a separable Borel metric domain. The value of $\eta>0$ is chosen below from raw moment bounds.

\begin{lemma}[Mass-preserving inward compander]\label{lem:masscompander}
For every $M\ge1$ the open simplex has a Borel cover with at most $M$ representatives satisfying
\begin{equation}\label{eq:compander}
 d_H(\pi,Q_M\pi)\le C_{d,\eta}M^{-1/d}w(\pi),\quad
 D(Q_M\pi)\le D(\pi),\quad
 (Q_M\pi)_i\ge\pi_i/n\quad\hbox{for all }i.
\end{equation}
The uniform preparation is a representative. Thus the fixed admissibility relation
\begin{equation}\label{eq:massrelation}
 \mathcal D(\pi)=\{\rho:D(\rho)\le D(\pi),\ \rho_i\ge\pi_i/n\text{ for every }i\}
\end{equation}
is respected by the static cover, independently of $M$.
\end{lemma}
\begin{proof}
Choose the least index $j$ of a largest coordinate of $\pi$ and set $v_i=\log(\pi_j/\pi_i)\ge0$, so $v_j=0$. For an integer $K\ge1$ define
\[
 \widehat v_i=-\eta^{-1}\log\left(\frac{\lceil K e^{-\eta v_i}\rceil}{K}\right),
 \qquad Q\pi_i=\frac{e^{-\widehat v_i}}{\sum_h e^{-\widehat v_h}}.
\]
There are at most $nK^d$ representatives: at least one deficit is zero and each other transformed coordinate takes $K$ values. All cells are Borel. For $u=e^{-\eta v_i}$ and $\widehat u=\lceil Ku\rceil/K$, one has $u\le\widehat u\le u+K^{-1}$, including the unbounded last cell. Consequently
\[
 0\le v_i-\widehat v_i\le e^{\eta v_i}/(\eta K).
\]
Since one error is zero, their oscillation is at most their maximum, bounded by $w(\pi)/(\eta K)$. This is exactly the projective error. Also $\max_i\widehat v_i\le\max_i v_i$, proving inwardness. Finally $\sum_i e^{-v_i}\ge1$ and $\sum_i e^{-\widehat v_i}\le n$, so $Q\pi_i/\pi_i\ge1/n$.

For $M\ge n$ choose $K=\lfloor(M/n)^{1/d}\rfloor\ge\tfrac12(M/n)^{1/d}$. For $M<n$ use the single uniform representative. Its error is $D(\pi)\le e^{\eta D(\pi)}/(e\eta)$, so enlarging the dimension-dependent constant proves \eqref{eq:compander} for all budgets. The preparation has all deficits zero.
\end{proof}

\begin{lemma}[Raw block certificate]\label{lem:hmmraw}
The raw experiment \eqref{eq:emissions}--\eqref{eq:bayes}, together with \eqref{eq:massrelation}, satisfies all kernel conditions of Theorem~\ref{thm:main}, with $L=1$ and explicit finite constants depending only on the fixed parameters in \eqref{eq:hmmparams}, $P,m,s$ and $d$. They are uniform in horizon and in the declared report-adapted action choices. Individual transitions need not have a strict projective contraction coefficient.
\end{lemma}
\begin{proof}
For a nonnegative matrix with no zero column, differentiate $\log\sum_i e^{z_i}Q_{ij}$ in a direction $h$. Every derivative is a convex combination of the $h_i$, so oscillation cannot increase after integrating along a log-coordinate segment. Multiplication by the same positive likelihood vector cancels in the projective difference. Thus every one-step update is nonexpansive, including the missing report.

Fix a cube $G=[-b,b]^d$ of positive volume. Uniform bounds on this cube are
\[
 g_-=q_-(2\pi\sigma^2)^{-d/2}
       e^{-d(b+R)^2/(2\sigma^2)},\qquad
 g_+=q_+(2\pi\sigma^2)^{-d/2}.
\]
Given any visible past, the probability of an observed report in $G$ is at least $p_0=g_-\lvert G\rvert>0$. The event $E$ of $m$ such reports has conditional probability at least $p=p_0^m$, without an independence assertion. On this event the matrix of the exact unnormalized block update,
$Q_{s+1}\operatorname{diag}(g(Y_{s+1}))\cdots
Q_{s+m}\operatorname{diag}(g(Y_{s+m}))$,
has every entry between $l=s^m g_-^m$ and $u=g_+^m$.

For any positive matrix with entries in $[l,u]$, its log-coordinate derivative probability vector for each output coordinate dominates $(l/u)$ times the common input probability vector. Subtract this common mass. The oscillation of the derivative is at most $(1-l/u)\osc(h)$, and integration gives contraction coefficient $\chi=1-l/u<1$. Off $E$ use nonexpansiveness. Therefore \eqref{eq:blockpair} holds with
\begin{equation}\label{eq:hmmkappa}
 \kappa=1-p(1-\chi^2)<1.
\end{equation}

It remains to control the \emph{rounded} trajectory inside that same block. Put $\beta=1/n$. On $E$, Bayes normalization is at most $g_+$, and every relation-preserving rounded state satisfies, coordinatewise,
\[
 z_j\ge\beta(g_-/g_+)z_{j-1}Q_{s+j}.
\]
Induction and primitive support imply
\begin{equation}\label{eq:floor}
 (z_m)_i\ge c_*:=\beta^m(g_-/g_+)^m s^m>0
 \quad\hbox{for every }i.
\end{equation}
Hence $D(z_m)\le R_*:=\log(1/c_*)$, independently of the candidate at the block start. This is precisely where componentwise mass preservation, beyond inwardness, is used.

For an arbitrary report put $V_a(y)=\osc_i\log g_i^a(y)$. The column sums of $Q$ lie in $[s,n]$, so
\[
 D(F_t(z,a,y))\le D(z)+C_s+V_a(y),\qquad C_s=\log(n/s).
\]
Relation-preserving rounding cannot increase this spread. The Gaussian likelihood ratios are affine in $y$ and the missing likelihood ratios are bounded. Thus for any fixed $\eta_0>0$,
\[
 \mathcal M=\sup_{x,a}\int e^{2\eta_0 V_a(y)}J_a(dy\mid x)<\infty.
\]
For completeness $V_a(y)\le v_0+v_1\norm{y}_1$ on the continuous part, with $v_1=2R/\sigma^2$ and $v_0$ at least both
$\log(q_+/q_-)+dR^2/(2\sigma^2)$ and
$\log((1-q_-)/(1-q_+))$.
The Gaussian bound $\E e^{c\norm{Y}_1}\le2^d e^{cdR+dc^2\sigma^2/2}$ supplies an explicit finite upper for $\mathcal M$, uniformly over hidden-state mixtures.

Let $Z=\sum_{j=1}^m(C_s+V_{a_{s+j}}(Y_{s+j}))\ge0$. Conditional iteration gives
$\E e^{2\eta_0 Z}\le C_0:=(e^{2\eta_0 C_s}\mathcal M)^m$.
Choose $0<\eta\le\eta_0$ with
$(\eta/\eta_0)(C_0-1)\le p/2$.
Convexity yields $e^{2\eta Z}-1\le(\eta/\eta_0)(e^{2\eta_0 Z}-1)$, whence
\[
 \E[e^{2\eta Z}\one_{E^c}\mid\cF_s]\le1-p/2=:a<1.
\]
On $E$ use \eqref{eq:floor}; on $E^c$ use $D(z_m)\le D(z_0)+Z$. This proves \eqref{eq:robustdrift} with $b=e^{2\eta R_*}$. The same growth estimate for each prefix proves \eqref{eq:within} with $A=C_0$, $B=0$. All estimates hold under the true report law for arbitrary adapted selections in \eqref{eq:massrelation}. No invariant distribution, observation window, or approximate-law drift is used.
\end{proof}

\begin{theorem}[An acquired law for sparse Gaussian hidden-state experiments]\label{thm:hmm}
For the experiment above fix $T\ge m$, $M\ge1$ and $b_{\rm out}\ge1$. Require each terminal forecast to lie in $\{j2^{-b_{\rm out}}:0\le j\le2^{b_{\rm out}}\}$. Set
\[
 A_T=\frac1{16d}\sum_{i=1}^d\E[\pi_{T,i}(1-\pi_{T,i})],\qquad
 b_\delta=\frac14-\frac1{16d}\sum_{i=1}^d
 \overline\pi_{T,i}(1-\overline\pi_{T,i})-\delta^2.
\]
For constants $0<c<C<\infty$ depending only on the fixed class parameters,
\begin{equation}\label{eq:hmmlaw}
 c\{A_T+M^{-2/d}+2^{-2b_{\rm out}}+\delta^2\}
 \le\mathcal R_\cp-b_\delta
 \le\mathcal R_\on-b_\delta
 \le C\{A_T+M^{-2/d}+2^{-2b_{\rm out}}+\delta^2\}.
\end{equation}
The same matched memory, output and calibration bounds hold after subtracting the common full-history acquisition term $A_T$. The risks are sign-minimax risks for the one indexed task \eqref{eq:filtertask}. All constants are uniform over the certified exploration rules for fixed $d$ and fixed parameter bounds; controller states are additional resources. The exponent $d$ is certified by positive actual acquired mass.
\end{theorem}
\begin{proof}
The physical-oracle variance of \eqref{eq:filtertask}, averaged over the probe, is $b_\delta$ for either sign. The acquisition uncertainty is $A_T$, by conditional projection of each indicator of $Z_T$. The score map is $1/4$-Lipschitz from $d_H$ to $\norm{\cdot}_d$: differentiation of a softmax coordinate gives $\pi_i(h_i-\sum_j\pi_jh_j)$, bounded in absolute value by $\osc(h)$. Lemmas~\ref{lem:masscompander} and \ref{lem:hmmraw} and Theorem~\ref{thm:main} yield nominal online score distortion $CM^{-2/d}$, from the exact uniform seed. Dyadic rounding adds at most $2^{-b_{\rm out}}$ in root task error, and Proposition~\ref{prop:calibration} gives the upper.

For the lower retain only the actual event $E_T$ that the last $m$ reports are observed and lie in $G$. Its probability is at least $p$. Conditional on a history at time $T-m$ and the first $m-1$ reports in this event, the prediction $r=\pi_{T-1}Q_T$ has every coordinate at least
\[
 c_r=(g_-/g_+)^{m-1}s^m.
\]
This follows by the unrounded coordinatewise argument used for \eqref{eq:floor}. The last action is already selected before $Y_T$. For that fixed action the map $y\mapsto(\pi_{T,1},\ldots,\pi_{T,d})$ is injective: its log ratios equal $B_{a_T}y$ plus a known vector. Its Jacobian determinant is
\begin{equation}\label{eq:gaussjac}
 \abs{\det B_{a_T}}\prod_{i=0}^d\pi_{T,i}
 \ge b_0 c_p^{d+1},\qquad c_p=c_r(g_-/g_+).
\end{equation}
The conditional continuous report density is at most $g_+$. Thus, restricting to $y\in G$ and scaling the task coordinates by $1/4$, the resulting score submeasure has Lebesgue density at most
\begin{equation}\label{eq:actualdensity}
 C_{\rm dens}=4^d g_+/(b_0c_p^{d+1}).
\end{equation}
Integrating over the unrestricted distribution of the conditioning prefix, with its event indicator, preserves this upper bound. No smooth density for the remaining posterior law is required. The actual score submeasure on $E_T$ has mass at least $p$ and small balls of $\norm{\cdot}_d$ have mass at most $C_{\rm dens}v_dd^{d/2}r^d$. Proposition~\ref{prop:acquired} therefore gives $q_J(\nu_T)\ge cJ^{-2/d}$ for every $J\ge1$, with the constant reduced for any finitely many small budgets if necessary; the bounded-density ball bound itself is valid at every radius.

A retained label determines one vector of possible probe forecasts, drawn from a grid of size $(2^{b_{\rm out}}+1)^d$. Averaging the two sign risks and applying Proposition~\ref{prop:calibration} gives the lower
\[
 A_T+\delta^2+c\min\{M,(2^{b_{\rm out}}+1)^d\}^{-2/d}.
\]
The last power equals
$\max\{M^{-2/d},(2^{b_{\rm out}}+1)^{-2}\}$.
Since $2^{b_{\rm out}}+1\le2^{b_{\rm out}+1}$ and the maximum of two nonnegative numbers is between half their sum and their sum, this is comparable to $M^{-2/d}+2^{-2b_{\rm out}}$. This proves \eqref{eq:hmmlaw}, and proves the stated bound on excess after subtracting $A_T$ rather than merely comparing the larger total risks.
\end{proof}

For example the four-state support of a lazy directed cycle, with positive entries only on each self-loop and its successor edge, is primitive but has one-step projective coefficient one because some row supports are disjoint. Taking any such transition weights bounded below, affinely independent Gaussian means, and unequal $q_i$ gives a nonempty class to which the block argument genuinely adds something. It does not manufacture a strict one-step coefficient by changing a reported missing value into an independent hold.

\subsection{Certified finite numerical interfaces}
The construction above is an exact Borel machine. Here is a finite enclosure route which preserves its load-bearing relation. Given certified intervals $[l_i,u_i]$ for unnormalized posterior log weights, each of width at most $h$, define
\[
 d_i=\max\{0,\max_jl_j-u_i\}.
\]
For the true deficits $v_i=\max_j z_j-z_i$, one has $0\le d_i\le v_i$ and $v_i-d_i\le2h$. At least one $d_i$ is zero, and every true maximal coordinate has $d_i=0$. Apply the deficit compander to $d_i$. Its output still has at most $nK^d$ possibilities, is inward, and dominates the true candidate coordinatewise by $1/n$. Its projective error is at most $r_M w+2h$. Thus \eqref{eq:localerror} holds with $u=2h$ without guessing a maximizing coordinate or assuming exact comparison at a cell boundary.

Known computable calibration parameters and interval routines for sums, exponentials and logarithms give finite such enclosures on any bounded input box. Because Gaussian reports are unbounded, a uniform finite packet cannot encode every exact report. With a declared box $[-U,U]^d$, $U>R$, the probability of any out-of-box continuous report before time $T$ is at most
\[
 2dT\exp(-(U-R)^2/(2\sigma^2)).
\]
Charge an explicit failure marker and use the bounded-loss cost of that event; inside the box use the enclosure construction. Input bins, workspace, program constants and arithmetic time are thereby finite and separately recorded. This finite-interface comparison uses the analogue task benchmark plus its certified input defect; it does not identify the digitally acquired posterior law with the analogue law used in \eqref{eq:actualdensity}. A digital-only checkpoint identity instead uses its own actual bin-history law. No horizon-independent finite packet claim follows from an exact Borel input.

\section{Finite acquisition through singular nonreset dynamics}\label{sec:singular}
The next verification has a finite atomic acquired law, a singular physical law, expanding updates and changing mode masses. It provides a second realization of the general theorem, not a premise of that theorem.

Let $r_n\in[1/5,1/3]$, $\ell_0=1$, $\ell_n=\prod_{j=1}^n r_j$, and
\[
 x(\omega)=\sum_{n\ge1}(1-r_n)\ell_{n-1}\omega_n,
 \qquad \omega\in\{0,1\}^{\N}.
\]
The physical Cantor set $K$ is the image of this map. For a word $u$ of length $h$, let
\[
 c(u)=\sum_{n=1}^h(1-r_n)\ell_{n-1}u_n+\tfrac12\ell_h.
\]
Adjoin all these fair-tail means to $K$ to obtain $K^*$. It is compact: centers of unbounded depth approach $K$, and bounded depths give only finitely many centers. Finite centers represent posterior distributions, not additional physical Cantor points.

Prepare a mode $I_0\in\{0,1\}$ with any declared distribution, independently of a fair infinite tail. Report the mode and make $B$ paid calls revealing successive initial bits. At each of $T$ further opportunities, independently activate with any deterministic probability $\vartheta_t$. An inactive call copies both physical variables and costs one call. On activation, use the following raw stochastic operations, with $\gamma_0=1/4096$ and arbitrary deterministic $\alpha_t\in[0,1]$:
\begin{center}\small
\begin{tabular}{cccl}
\toprule
Source & Destination & Probability & Tail operation\\\midrule
$0$&$0$&$(1-\alpha_t)(1-\gamma_0)$&prepend two fair bits\\
$0$&$0$&$(1-\alpha_t)\gamma_0$&delete first bit\\
$0$&$1$&$\alpha_t$&delete first bit\\
$1$&$0$&$1/2$&prepend six fair bits\\
$1$&$1$&$1/2$&copy\\\bottomrule
\end{tabular}\end{center}
Report destination, operation and fresh prepended bits, but neither the deleted bit nor the old surviving tail. All kernels are nonnegative and normalized; operation probabilities are independent of bit values. Both mode changes transport a nonconstant old tail.

Let $w_T=\PP(I_T=1)$ under this actual protocol. The terminal Bernoulli probe has mean $X+s\delta$, where
\begin{equation}\label{eq:singulartask}
 X=\tfrac12+\tfrac14(I_T-w_T)+\tfrac18(x(\omega_T)-\tfrac12),
 \qquad s\in\{-1,1\},\quad0\le\delta\le1/16.
\end{equation}
These are valid means, and $\E X=1/2$. The mode and acquired cylinder posterior determine every legal continuation. Their score intervals are disjoint across modes and the current probe is strictly affine in the cylinder mean within a mode. They therefore realize the predictive quotient directly, including the declared time interface.

\begin{lemma}[Completed-domain geometry and raw gains]\label{lem:singularraw}
On $K^*$, deletion, fixing $c(\varnothing)$, is $30$-Lipschitz; prepending a fixed word of length $s$ is $6\,3^{-s}$-Lipschitz. On $S=\{0,1\}\times K^*$ put
\[
 d((i,x),(i,y))=\sqrt{v_i}|x-y|,
 \quad (v_0,v_1)=(1800,1),
 \quad d((0,x),(1,y))=2\sqrt{1800}.
\]
All centers of depth at most $h$ in both modes form an inward cover with
\[
 M_h=2^{h+2}-2,\qquad r_{M_h}=\sqrt{1800}\ell_h,\qquad w=1.
\]
The active true-report kernel satisfies the one-step pair inequality with $\kappa=17/20$, and its candidate envelope is identically one, independently of the mode law and schedule.
\end{lemma}
\begin{proof}
If two completed coordinates first differ, or one terminates, at level $k$, their separation lies between $\ell_{k-1}/6$ and $\ell_{k-1}$. Opposite branches are separated by $(1-2r_k)\ell_{k-1}\ge\ell_{k-1}/3$; a parent center and either branch are separated by $(1/2-r_k)\ell_{k-1}\ge\ell_{k-1}/6$. This also proves unique decoding of finite centers and infinite points. Deletion lowers the distinguishing level by one, giving gain at most $6/r_{k-1}\le30$. If the first digits differ or one point is the empty center, the gain is at most six. A common prepended word raises the distinguishing level by $s$, giving gain at most $6\ell_{k+s-1}/\ell_{k-1}\le6\,3^{-s}$.

The displayed metric satisfies the triangle inequality because within-mode diameters are at most $\sqrt{1800}$. Leave shorter centers fixed and truncate longer words to depth $h$. The cells are Borel, the error is at most $\sqrt{1800}\ell_h$, and the count includes both the mode and every possible retained word length.

For a candidate in the wrong mode, apply the reported operation to its coordinate and set its mode to the reported destination. This defines the common Borel update on every true/candidate input. For same-mode pairs the matrix of squared gains, with rows indexed by source and columns by destination, is
\[
 C_t=\begin{pmatrix}
 (1-\alpha_t)85/128&900\alpha_t\\
 2/59049&1/2
 \end{pmatrix}.
\]
Indeed $(1-1/4096)(4/9)+900/4096=85/128$. For $v=(1800,1)^{\mathsf T}$ both rows obey $(C_tv)_i\le(17/20)v_i$: the first is maximized at the larger of $1800(85/128)$ and $900$, and the second equals $3600/59049+1/2<17/20$. For cross-mode pairs the squared distance before the update is $4\cdot1800$ and afterward at most $1800$. The gain is at most $1/4$. These are inequalities under the physical report law for all pairs, not just a realized code trajectory. The constant envelope verifies the relation-robust conditions with $a=A=0$, $b=B=1$.
\end{proof}

\begin{theorem}[Actual-depth acquisition--memory law]\label{thm:singular}
Let $H_0=B$, and update the acquired depth by $h+2$, $h+6$, $\max(h-1,0)$ or $h$ according to the reported operation. Set $\omega_{T,i,h}=\PP(I_T=i,H_T=h)$. The actual nominal score law is
\begin{equation}\label{eq:atomic}
 \nu_{B,T}=\sum_{i,h}\omega_{T,i,h}2^{-h}
   \sum_{u\in\{0,1\}^h}
    \delta_{\,1/2+(i-w_T)/4+(c(u)-1/2)/8}.
\end{equation}
It is finite atomic. Put $A_{B,T}=\E(X-\E[X\mid\cH_{B,T}])^2$ and $b_\delta=\E[X(1-X)]-\delta^2$. For $M\ge6$ and $n=\lfloor\log_2 M\rfloor$, the sign-minimax risks satisfy
\begin{align}
 \mathcal R_\cp-b_\delta&=\delta^2+A_{B,T}+q_M(\nu_{B,T}),\label{eq:singularexact}\\
 \mathcal R_\cp-b_\delta&\asymp\mathcal R_\on-b_\delta
 \asymp\delta^2+\E\ell_{\min(H_T,n)}^2.\label{eq:singularlaw}
\end{align}
The constants are uniform over $B,T$, all ratios and schedules, and all initial mode laws. Raw acquisition is exactly $B+T+2$, including preparation, every inactive opportunity, and the terminal probe.
\end{theorem}
\begin{proof}
Inductively the full history knows a finite prefix, the unrevealed remainder is fair, and the prefix is uniform conditional on its depth and mode. Prepending discloses fresh independent bits, deletion reveals no unknown bit, and holds change nothing. This proves \eqref{eq:atomic}, with $H_T\le B+6T$. The unresolved physical coordinate variance in a depth-$h$ cylinder lies between $\ell_h^2/9$ and $\ell_h^2/4$: the first unresolved bit alone contributes at least $\ell_h^2/9$, and the whole cylinder has length $\ell_h$. Hence
\begin{equation}\label{eq:singularacquisition}
 \frac1{576}\E\ell_{H_T}^2\le A_{B,T}\le\frac1{256}\E\ell_{H_T}^2.
\end{equation}

Choose $h=\lfloor\log_2(M+2)\rfloor-2\ge1$. During the initial paid readings retain only the first $\min(B,h)$ bits. The seed is legal, has envelope one and error at most $\sqrt{1800}\ell_h$. Subsequently prepend and truncate, delete the first retained bit if present, or copy. Lemma~\ref{lem:singularraw}, Theorem~\ref{thm:main} with $m=1$, and Corollary~\ref{cor:holds} give nominal acquired-score distortion at most $C\ell_h^2\le C'\ell_n^2$, since $0\le n-h\le2$. The score is Lipschitz also across modes. The retained word is always a genuine prefix. Its origin pattern depends on the operations, not bit values; conditioning on that pattern and then averaging shows that the center readout is the conditional mean of $X$ given its retained word and mode. It is unbiased, so \eqref{eq:sign} adds exactly $\delta^2$.

For a converse, every operation preserves a fair physical tail conditional on the terminal mode. At least one terminal mode has mass at least $1/2$. Level-$k$ physical cylinders have length $\ell_k$ and gaps at least $\ell_k$. Their score gaps are at least $\ell_k/8$. A score ball of radius $\ell_k/32$ has diameter $\ell_k/16$, so meets at most one such cylinder of the selected mode. Its conditional mass is $2^{-k}$. Taking $k=n+2$ gives $2^{-k}\le1/(2M)$ and $\ell_k\ge\ell_n/25$. The union argument yields $q_M(\law(X))\ge c\ell_n^2$. This lower witness is the physical-oracle law; no small-ball regularity of the actual atomic law is asserted.

For every checkpoint forecast, averaging the two signs gives
\[
 \sup_s R_s-b_\delta\ge\delta^2+A_{B,T}+q_M(\nu_{B,T}).
\]
The finite atomic law has an optimal nominal code; replace its readouts by label-conditional means. This preserves its label count and makes its mean error zero, so
\[
 \sup_s R_s-b_\delta=\delta^2+A_{B,T}+q_M(\nu_{B,T})
\]
for such an optimal code. This proves \eqref{eq:singularexact}. Proposition~\ref{prop:acquired} and \eqref{eq:singularacquisition} compare both risks with $\delta^2+\E\ell_{H_T}^2+\ell_n^2$. Finally
\[
 \max\{\E\ell_{H_T}^2,\ell_n^2\}
 \le\E\ell_{\min(H_T,n)}^2
 \le\E\ell_{H_T}^2+\ell_n^2,
\]
which proves \eqref{eq:singularlaw}.
\end{proof}

This realization includes expanding deletion, arbitrary nonstationary mode masses and a nonhomogeneous sequence of singular scales. Its finite actual support changes with acquisition. The matching law is not obtained by assigning the physical Cantor dimension to the atomic acquired distribution. The known ratios may be supplied by a certified computable routine or a charged finite table to the required depth; arbitrary real constants do not by themselves furnish a finite program. The exact prefix update needs only the counted finite word, whereas numeric evaluation of the final center needs its own enclosure and program account.

\section{Stopped causal morphisms and common-risk composition}\label{sec:morphisms}
We now fix the measurable and computational meaning of experiment comparison. The direction $\mathfrak E\rightsquigarrow\mathfrak F$ means that the source $\mathfrak E$ implements the target $\mathfrak F$. It is not a comparison of two separately optimized controllers on a shared report stream.

\subsection{Measurable coupling and report continuity}
\begin{lemma}[A jointly measurable maximal coupling]\label{lem:coupling}
Let $P_z,Q_z$ be probability kernels from a standard Borel parameter space $Z$ to a standard Borel report space $Y$. There is a probability kernel $C_z$ on $Y\times Y$ with these marginals and
\[
 C_z\{(y,y'):y\ne y'\}=\TV(P_z,Q_z),
 \qquad \TV(P,Q):=\sup_A|P(A)-Q(A)|.
\]
\end{lemma}
\begin{proof}
Let $R_z=(P_z+Q_z)/2$. Choose increasing finite partitions of $Y$ whose union generates its Borel sigma field and separates points. Such partitions follow from a countable separating generator of a standard Borel space. On an atom $A$ put $p_k(z,y)=P_z(A)/R_z(A)$, with value zero when the denominator vanishes. These functions are jointly measurable and bounded by two. For every fixed $z$, the martingale convergence theorem applied to the Radon--Nikodym derivative $dP_z/dR_z$ gives a limit equal to that derivative $R_z$-almost everywhere. Define $p$ to be the clipped limsup and $q=2-p$. These are jointly measurable versions of the two densities for each $z$; their exceptional sets need not be independent of $z$.

The kernel $\mu_z(dy)=\min(p(z,y),q(z,y))R_z(dy)$ has measurable mass $c_z$. Its two residual measures have mass $1-c_z$ and are mutually singular. Set
\[
 C_z=\operatorname{diag}_*\mu_z+
 \frac{(P_z-\mu_z)\otimes(Q_z-\mu_z)}{1-c_z}
 \quad(c_z<1),
\]
and set $C_z=\operatorname{diag}_*P_z$ if $c_z=1$. Products of finite kernels and the displayed measurable case distinction are kernels. The diagonal is Borel. The residual product gives it zero mass by mutual singularity, so disagreement has probability $1-c_z=\TV(P_z,Q_z)$. This also proves the parameterized measurability needed for successive causal couplings.
\end{proof}

\begin{proposition}[Stopped report-law comparison]\label{prop:report}
Suppose the true and representative report kernels obey
$\TV(J_t(\cdot\mid x,a),J_t(\cdot\mid z,a))\le L_Jd(x,z)$
uniformly in the declared actions. Suppose the terminal task-report kernel has the analogous constant $L_Y$. Use the same declared controller, controller randomization and visible stopping rule $\tau\le T$ in the exact and simulated experiments. Let $\widehat S_t$ be the shadow recursion driven by the physical reports. Then their complete stopped report-and-task laws differ in total variation by at most
\begin{equation}\label{eq:stoppedtv}
 \min\left\{1,
 L_J\E\sum_{t=1}^{\tau}d(S_{t-1},\widehat S_{t-1})
 +L_Y\E d(S_\tau,\widehat S_\tau)\right\}.
\end{equation}
At a deterministic $T$, Theorem~\ref{thm:main} bounds this by
$\min\{1,(TL_J+L_Y)\mathcal E_M\}$ in the exact-readout case.
\end{proposition}
\begin{proof}
Use Lemma~\ref{lem:coupling} successively until the first discrepancy. Before that discrepancy the controllers see the same reports, choose the same actions and make the same stop decisions. Conditional disagreement probability at a continuing step is bounded by $L_Jd(S_{t-1},\widehat S_{t-1})$. Continue a physical-law shadow after any discrepancy, so its error estimates still apply without conditioning its law on successful coupling. Sum first-discrepancy probabilities and couple the final task on the agreement event. The indicators of continuing are the physical $\{t\le\tau\}$ up to disagreement; dropping the agreement indicator gives \eqref{eq:stoppedtv}. At deterministic times use Cauchy--Schwarz and the uniform root error. This proof does not replace a stopped sum by $\E\tau$ times an unconditional deterministic-time bound.
\end{proof}

For the Gaussian filter, report laws are mixtures of the same emission kernels with weights $\pi Q$. Total variation is at most $\tfrac12\norm{\pi Q-\rho Q}_1\le\tfrac12 d_H(\pi,\rho)$; the last inequality follows by differentiating the softmax along a log-coordinate segment. The indexed terminal probe has $L_Y=1/4$. In the singular experiment the report law does not depend on the coordinate within a mode; cross-mode total variation is at most one and the cross-mode distance is $2\sqrt{1800}$. Its affine probe is also Lipschitz. These are direct raw-kernel continuity checks. Uniform prediction risk does not imply infinite-transcript or large-deviation equivalence.

\subsection{Typed experiment morphisms}
Fix a common parameter set $\Lambda$ and two prepared families $\mathfrak E=(K^{E,\lambda},\mu^{E,\lambda})_{\lambda\in\Lambda}$ and $\mathfrak F=(K^{F,\lambda},\mu^{F,\lambda})_{\lambda\in\Lambda}$. All state, action, report, cost and terminal-outcome spaces are standard Borel. Source and target visible histories are finite products
\[
 H^E_n=Y^E_0\times\prod_{j=1}^n(A^E_j\times Y^E_j\times\mathcal C^E_j),
 \qquad H^F_i=Y^F_0\times\prod_{j=1}^i(A^F_j\times Y^F_j\times\mathcal C^F_j).
\]
Disjoint unions over lengths carry their usual standard Borel structure. Histories include a preparation report, all failures and publicly available resource records. Legal action correspondences are declared measurable subsets of $H\times A$. The specified strategy classes consist of parameter-independent Borel action kernels supported on these correspondences, together with bounded visible stop rules and terminal decision rules. All such rules may use permitted independent randomization, but not the unknown $\lambda$.

\begin{definition}[A certified stopped causal morphism]\label{def:morphism}
For a fixed target horizon $T$, a morphism $\mathfrak E\rightsquigarrow\mathfrak F$ consists of the following data and requirements.

\emph{Causal realization.} A parameter-independent transducer has a declared internal state space $\mathcal U$, initialization kernel, source-action kernels, report transformations, and state-update kernels. Each physical operation depends only on its current retained registers, the current virtual action, the current source report and its allowed fresh randomness. Equivalently its virtual history maps are causal on source stopped histories. A strategy lift $\mathscr L:\Pi_F\to\Pi_E$ specifies the resulting source policy for each declared target policy. It must be induced by those kernels, not by access to unretained histories. Its source actions are legal at every declared reachable interface.

\emph{Clock and stopping compatibility.} Virtual completion times $0\le\sigma_0\le\cdots\le\sigma_T$ are source stopping times, bounded by a declared physical-call budget $N(T)$. Virtual reports and costs at step $i$ are measurable at $\sigma_i$. For every target stop $\tau\le T$, $\sigma_\tau$ is a legal source stop; no subsequent source information is used in the virtual terminal decision. A phase/clock register and the cost transformation specify which physical actions realize each virtual action, including preparation, waits and failed attempts. Action availability after each transformed prefix must agree with the declared target interface.

\emph{Common task.} A parameter-independent Borel terminal map $\chi$ takes the actual source terminal outcome and available stopped registers to the common scored-outcome space $\mathcal O$. The target has its declared outcome in that same space. The decision map sends the retained tuple to the same decision space $\mathcal A_*$. A loss $\ell:\mathcal O\times\mathcal A_*\to[0,L_*]$ is fixed. Merely matching unscored transcript marginals is not a common-task certificate.

\emph{Uniform defect.} For every $\lambda\in\Lambda$, $\beta\in\Pi_F$ and its declared $\tau\le T$, let $\widetilde P^{\beta,\tau}_\lambda$ be the joint law of the virtual stopped history, permitted public policy coins, and transformed common outcome. Let $P^{F,\beta,\tau}_\lambda$ be the corresponding true target law. Require
\begin{equation}\label{eq:defect}
 \sup_{\lambda,\beta,\tau}
 \TV(\widetilde P^{\beta,\tau}_\lambda,P^{F,\beta,\tau}_\lambda)
 \le\varepsilon.
\end{equation}
All laws and comparisons use these same declared classes; a single-controller calculation certifies only the singleton class unless uniformity is proved.

\emph{Resource realization.} A finite implementation separately bounds predictor labels $M$, transducer states $R$, controller states $D$, phase/clock states $Q$, temporary writable bits $W$, read-only program bits $L_{\rm prog}$, calibration and input precision, raw calls and physical time. Any register consulted later is persistent and belongs to the counted tuple. Workspace must be erased before the next declared persistent cut or included there. Uniformly finite resource bounds are extra data, not implications of the Borel description.
\end{definition}

The definition includes the preparation stage $\sigma_0$ and permits zero-call virtual deterministic operations, provided the finite virtual horizon and clock rule prevent an uncharged infinite execution. A raw acquisition count does not count only successful virtual reports. Parameters in read-only constants must be known calibration, not unknown true values. Complete defect is stronger than one-step marginal closeness.

\begin{theorem}[Common-risk composition and a reverse converse]\label{thm:morphism}
A morphism in Definition~\ref{def:morphism} implements every certified target predictor with a serial source state count
\begin{equation}\label{eq:tuple}
 K\le MRDQ
\end{equation}
and source risk at most $r_\lambda+L_*\varepsilon$ when its target risk is $r_\lambda$. For two composable morphisms, assume the first lifted policy, its stops and common-task maps lie in the second certified class. Their complete defects add, capped at one; serial state overheads multiply. Raw calls and physical time are sums over the actually executed physical subroutines; nested per-call upper bounds may therefore multiply. Simultaneously required workspace and stored program descriptions admit the corresponding additive upper accounts.

In particular Theorem~\ref{thm:main} yields the common-risk upper
\begin{equation}\label{eq:composedrisk}
 R_E(MRDQ)\le B_T+(L_f\mathcal E_{M,u}+v)^2+L_*\varepsilon,
\end{equation}
where $\mathcal E_{M,u}$ denotes \eqref{eq:constants} at numerical allowance $u$, with the target benchmark $B_T$ and certified acquisition budget. Root metric errors are combined before squaring; complete bounded-loss defect is added afterwards.

Suppose additionally a reverse morphism converts every source $K$-state procedure in the comparison class to a target procedure with at most $\Psi(K)$ states and defect at most $\varepsilon_-$. For the same fixed preparation, task and risk criterion,
\begin{equation}\label{eq:reverse}
 R_E(K)\ge R_F(\Psi(K))-L_*\varepsilon_-.
\end{equation}
Thus actual target acquired-law lower bounds transfer at the explicitly inflated budget. The same implication holds for minimax risks when both certificates are uniform in the same parameter set.
\end{theorem}
\begin{proof}
Retain the tuple of predictor, transducer, controller and clock states. Run the prescribed source subroutine, transform its report, update the target predictor and use the compatible task map. Causality, action compatibility and the stopping conditions make this a legal source procedure. Its tuple set has at most $MRDQ$ elements. The bounded-loss inequality for the convention in Lemma~\ref{lem:coupling} gives risk difference at most $L_*\varepsilon$ from \eqref{eq:defect}, also after the common decision kernel is applied.

For composition insert the intermediate joint stopped law. Apply the triangle inequality and contraction of total variation under the downstream parameter-independent causal transformation. The class-closure hypothesis is needed to apply the second defect certificate to the lifted adaptive policy. Clock compatibility identifies the composed stopping time and common task. Serial storage of both tuples gives the product bound; summing all actual calls and time, including failures and preparation, gives the cost accounts. These are constructive upper accounts, not claims that every register is indispensable.

For the reverse assertion apply the same upper argument to an arbitrary source procedure and its target realization. It gives $R_F(\Psi(K))\le r_E+L_*\varepsilon_-$; take the infimum over source procedures. For minimax risks take the supremum in $\lambda$ before the infimum, using parameter-uniform certificates. The target and source Bayes baselines need not agree; no subtraction of separately recomputed baselines is performed.
\end{proof}

\begin{proposition}[An executable interface converse]\label{prop:interface}
Before a memory cut an experiment reports a tag $J\in\{1,\ldots,D\}$, with positive probabilities $p_j$, and a score $U$ of conditional law $\nu_j$. After the cut no report reveals either. A declared continuation requires exact recovery of $J$ as well as a squared forecast of $U$. The optimal distortion with $K$ retained states is
\begin{equation}\label{eq:tagallocation}
 \inf_{\substack{m_j\ge1\ \sum_jm_j\le K}}
     \sum_jp_jq_{m_j}(\nu_j),
\end{equation}
with infeasibility if $K<D$. Independent randomization cannot improve this value. For uniform tags and uniform $d$-cube laws the value is of order $(D/K)^{2/d}$ for $K\ge D$.
\end{proposition}
\begin{proof}
Fix a shared seed outside the null set on which exact recovery can fail. A retained label cannot have positive probability for two distinct tags, because its tag readout would have to be two different point masses. The label sets therefore split into disjoint supports of sizes $m_j$. Conditional projection and nearest-center quantization give the lower in \eqref{eq:tagallocation}; disjoint codebooks attain its infimum arbitrarily closely. Condition on the independent seed to cover randomized schemes. Cube small balls give $q_m\ge c_dm^{-2/d}$ and a regular grid gives the reverse bound. Jensen's inequality and allocation of at least $\lfloor K/D\rfloor$ centers to each tag prove the stated order.
\end{proof}

An executable challenge to a tuple of independent interface tags can make its product cardinality necessary. Without such a challenge or a reverse certificate, $RDQ$ is only an upper overhead. Similarly a defect allowance does not force an error floor. Its linear bounded-loss scale is nevertheless attainable: if the target observes a fair bit $U$ exactly and the source erases it with probability $e$, guessing a fair virtual bit on erasure has joint-law defect $e/2$ from $(U,U)$. The source's optimal squared prediction risk for $U$ is $e/4$ (output $1/2$ on erasure), while the target risk is zero. Three retained labels suffice for the source. This proves order-sharpness for that specified erasure interface, not a universal positive lower for every experiment with an upper defect bound.

\section{Relation to earlier results and quantitative scope}\label{sec:comparison}
\subsection{The mechanism and its antecedents}
Conditional projection, finite-center quantization and small-ball converses are standard; see Graf and Luschgy~\cite{graf}. Predictive representations and future-equivalence have substantial antecedents in Littman, Sutton and Singh~\cite{psr}, Shalizi and Crutchfield~\cite{shalizi}, and kernel sufficiency in Fritz~\cite{fritz}. Proposition~\ref{prop:quotient} specifies the measurable-image and pointwise instrument conditions needed here; it does not claim to originate predictive sufficiency. Blackwell~\cite{blackwell}, Le Cam~\cite{lecam} and Torgersen~\cite{torgersen} supply the classical comparison-of-experiments background. Our stopped causal definition adds an explicit resource and implementation interface to the common-loss comparison; total-variation contraction itself is standard.

Average contraction of random maps is classical~\cite{diaconis}. A negative mean logarithmic gain alone is not the square-moment certificate in Theorem~\ref{thm:main}. The theorem's contribution is the following implication: a budget-independent relation preserved by a static cover, raw block stability and relation-robust physical-law envelope bounds construct a hard-label recursion at every raw call, while actual acquired mass supplies its checkpoint converse. Lemma~\ref{lem:masscompander} and the primitive-support verification show that this implication reaches sparse Gaussian filters, rather than merely repeating a one-step positive-matrix calculation. Neither the shadow process nor the observation block is a retained resource.

The earlier one-step inward result is recovered by $m=1$. The completed-domain singular calculation and exact finite-acquisition law are retained as a complementary realization. Earlier allocation-weighted, countable-stratum, reset-preserved and continuation-information results in the companion mathematical supplement~\cite{supplement} keep their own hypotheses. A single global relation-robust cover is not asserted to subsume every such allocation theorem, and none of those realization-specific results is used to prove Theorem~\ref{thm:main}.

\subsection{The closest operational comparisons}
Recursive filter quantization is not new. Pag\`es and Pham~\cite{pages} analyze discrete-observation filter approximation using optimal quantization, while Sellami~\cite{sellami} develops a first-order quantization construction. These results motivate error transport, but a signal grid or a numerical approximation order is not automatically a fixed cardinality for every retained predictive state. Here the codebook is finite on an unbounded stable coordinate, and a positive actual posterior submeasure supplies the matching converse. We do not claim novelty for Bayes recursion, projective nonexpansiveness, or Gaussian likelihood calculations separately.

Theorem~9 of Subramanian, Sinha, Seraj and Mahajan~\cite{ais} assumes an approximate information-state generator, with reward and transition discrepancies measured on all histories. It bounds approximate dynamic-program values through an accumulated error recursion, and the resulting policy loses at most twice that accumulated bound. This concerns policy optimization and permits a broad compressed state space. Our theorem instead constructs a specified hard finite alphabet from raw kernels and a cover, and derives acquired-law prediction lower bounds. Our same-controller comparison does not imply their control conclusion; their generator conditions alone do not imply our cardinality law.

Theorem~16 of Kara and Y\"uksel~\cite{kara}, under their Assumption~4 and discount restriction, bounds the loss of finite-window control policies exponentially in window length. Their construction assumes finite observation and action alphabets and known kernels. Our Gaussian report alphabet is continuous, the retained object is one of $M$ states rather than a window, and the objective is a specified prediction task with an actual geometric converse. Our constants may be poor near instability; no superiority in control cost or algorithmic time is inferred from the different memory parameter.

Ghomi, Linder and Y\"uksel~\cite{ghomi} prove existence of stationary deterministic zero-delay codes for linear vector Markov sources with quadratic average distortion, with finite-horizon performance approaching the infinite-horizon value at order $1/T$. Wood, Linder and Y\"uksel~\cite{wood} establish corresponding structural results for finite irreducible aperiodic Markov sources and extensions with noisy communication. Their optimized coding policies and encoder/decoder information states are not the hard predictive-quotient alphabet here. In particular their horizon approximation rate is not the $M$-center acquired-mass lower bound of this paper. Conversely we do not establish their average-cost optimal stationary coding conclusions.

Theorem~1 of Kara~\cite{kara-learning} proves convergence of a finite-memory linear policy-evaluation iteration under its ergodicity and step-size assumptions to a projected Bellman fixed point. Its further error and control analysis includes projection, filter stability and exploration conditions. Our known-kernel enclosure calculation is not a learning theorem, and no statistical estimation guarantee for unknown dynamics is extracted from it. Nonanticipative rate-distortion theory~\cite{nonanticipative} provides another causal perspective, but a mutual-information budget is not a pointwise bound on a retained alphabet without a separate conversion argument.

\subsection{Condition numbers and resource scope}
The uniformity statements always hold at fixed class parameters. The following quantities make their deterioration explicit.
\begin{center}\small
\begin{tabular}{p{.29\textwidth}p{.63\textwidth}}
\toprule
Quantity & Role and loss of uniformity\\\midrule
$1-\sqrt\kappa$, $1-a$ & Denominators in the error radius and envelope. No stability-boundary uniformity.\\
$m,L,\Lambda_m$ & Propagation within a block. Long blocks or expansion enlarge constants.\\
$p_0^m$, $s^m$, $g_-/g_+$ & Primitive-filter good-block mass and coordinate floor; may be small.\\
$\eta$, $C_0$, $n$ & Compander resolution and robust moment constants; dimension is fixed.\\
$b_0$, $c_p$, $C_{\rm dens}$ & Actual Jacobian and acquired small-ball constants; degenerating means weaken the lower.\\
$\zeta_T$, $e_0$ & Acquired mass and seed error remain explicit; neither may be silently normalized away.\\
$\E\lceil\tau/m\rceil$ & Stopped selection/acquisition cost; deterministic-time uniformity does not remove it.\\
$R,D,Q,\varepsilon$ & Serial state overhead and complete common-task defect, not a free simulator.\\\bottomrule
\end{tabular}\end{center}

The general upper combines errors in one task norm and the geometric converses give matching laws for the specified memory, acquisition, output and calibration interfaces. The reverse-morphism theorem and executable tag challenge provide additional resource converses when their stated data are available. A necessary-and-sufficient characterization of recursive attainability of every acquired geometry is not proved. Nor is the full optimal region for temporary workspace, program length, arithmetic time, minimal simulator state, arbitrary adaptive acquisition design, and unknown-kernel learning. The bounded stopping theorem proves a new selection-sensitive estimate, not a horizon-free optimal stopping law. The complete-defect comparison is stopped and finite, not an infinite-path or large-deviation assertion.

These distinctions leave the central problem intact: experiments generate predictive objects; admissible finite recursions must realize their acquired geometry; and common-task causal morphisms must pay for their implementation. The theorem provides that connection over the displayed general class, with two structurally distinct raw-kernel verifications and explicit boundaries on what the connection does not yet characterize.

\begin{thebibliography}{99}
\bibitem{blackwell} D. Blackwell, \emph{Equivalent comparisons of experiments}, Ann. Math. Statist. \textbf{24} (1953), 265--272.
\bibitem{diaconis} P. Diaconis and D. Freedman, \emph{Iterated random functions}, SIAM Rev. \textbf{41} (1999), 45--76.
\bibitem{fritz} T. Fritz, \emph{A synthetic approach to Markov kernels, conditional independence and theorems on sufficient statistics}, Adv. Math. \textbf{370} (2020), 107239.
\bibitem{ghomi} M. Ghomi, T. Linder and S. Y\"uksel, \emph{Zero-delay lossy coding of linear vector Markov sources: optimality of stationary codes and near optimality of finite memory codes}, IEEE Trans. Inform. Theory \textbf{68} (2022), 3474--3488. arXiv:2103.10810v3.
\bibitem{graf} S. Graf and H. Luschgy, \emph{Foundations of Quantization for Probability Distributions}, Lecture Notes in Mathematics 1730, Springer, 2000.
\bibitem{kara} A. D. Kara and S. Y\"uksel, \emph{Near optimality of finite memory feedback policies in partially observed Markov decision processes}, J. Mach. Learn. Res. \textbf{23} (2022), no.~11, 1--46.
\bibitem{kara-learning} A. D. Kara, \emph{Learning POMDPs with linear function approximation and finite memory}, arXiv:2505.14879v1, 2025.
\bibitem{lecam} L. Le Cam, \emph{Asymptotic Methods in Statistical Decision Theory}, Springer, New York, 1986.
\bibitem{psr} M. L. Littman, R. S. Sutton and S. Singh, \emph{Predictive representations of state}, Advances in Neural Information Processing Systems \textbf{14} (2001), 1555--1561.
\bibitem{pages} G. Pag\`es and H. Pham, \emph{Optimal quantization methods for nonlinear filtering with discrete-time observations}, Bernoulli \textbf{11} (2005), 893--932.
\bibitem{supplement} Q. Qi, \emph{General Theta Foundations I: Acquired Geometry and Causal Resource Transfer}, companion mathematical supplements, 2026.
\bibitem{sellami} A. Sellami, \emph{Quantization based filtering method using first order approximation}, SIAM J. Numer. Anal. \textbf{47} (2010), 4711--4734.
\bibitem{shalizi} C. R. Shalizi and J. P. Crutchfield, \emph{Computational mechanics: pattern and prediction, structure and simplicity}, J. Stat. Phys. \textbf{104} (2001), 817--879.
\bibitem{nonanticipative} C. D. Charalambous, P. A. Stavrou and N. U. Ahmed, \emph{Nonanticipative rate distortion function and relations to filtering theory}, IEEE Trans. Automat. Control \textbf{59} (2014), 937--952.
\bibitem{ais} J. Subramanian, A. Sinha, R. Seraj and A. Mahajan, \emph{Approximate information state for approximate planning and reinforcement learning in partially observed systems}, J. Mach. Learn. Res. \textbf{23} (2022), no.~12, 1--83.
\bibitem{torgersen} E. Torgersen, \emph{Comparison of Statistical Experiments}, Cambridge University Press, 1991.
\bibitem{wood} R. G. Wood, T. Linder and S. Y\"uksel, \emph{Optimal zero delay coding of Markov sources: stationary and finite memory codes}, IEEE Trans. Inform. Theory \textbf{63} (2017), 5968--5980; arXiv:1606.09135v2 consulted.
\end{thebibliography}

\end{document}
